% !TeX root = ../tfg.tex
% !TeX encoding = utf8
%
% ====================================================================
%  CAPÍTULO 3 — TEMA 3: PLANIFICACIÓN DE RECURSOS HUMANOS
%
%  Este capítulo cubre todos los apartados del Tema 3 de la asignatura:
%    3.1  Planificación de RRHH o de plantilla
%    3.2  Análisis de oferta y demanda de RRHH
%    3.3  Equilibrio entre oferta y demanda
%    3.4  Criterios de evaluación y sistemas de información (SIRH)
%    3.5  Procesos sustractivos
%    3.6  Alternativas a la reducción / incorporación de plantilla
% ====================================================================

\chapter{Tema 3: Planificación de Recursos Humanos}


% ================================================================
\section{Planificación de recursos humanos en CaixaBank}
% ================================================================

% RECORDATORIO AL EQUIPO:
% La planificación de RRHH consiste en anticipar las necesidades
% futuras de personal y definir las acciones necesarias para satisfacerlas.
%
% NIVELES DE PLANIFICACIÓN:
%   1. OPERATIVA (corto plazo, 1 año): cubrir necesidades inmediatas.
%   2. TÁCTICA (medio plazo, 1-3 años): ajuste de plantilla, formación,
%      movilidad interna.
%   3. ESTRATÉGICA (largo plazo, +3 años): alineada con la estrategia
%      de negocio; transforma la estructura de la plantilla para los
%      retos futuros (digitalización, nuevos mercados, etc.).
%
% PARA ESTA SECCIÓN INDICAD:
%   - ¿Existe un proceso formal de planificación de plantilla en CaixaBank?
%   - ¿Es operativa, táctica o estratégica (o las tres)?
%   - ¿Está integrada con el plan estratégico del grupo?
%   - ¿Está la empresa satisfecha? ¿Qué mejoraría?

La planificación de recursos humanos es el proceso mediante el cual la
organización prevé sus necesidades futuras de personal y define las
acciones necesarias ---contratación, formación, reducción,
reconversión--- para alinear la dotación de personal con los objetivos
estratégicos del negocio.

\subsection{¿Existe planificación formal de RRHH en CaixaBank?}

% TODO: CaixaBank publica planes estratégicos trienales (p.ej. Plan
% Estratégico 2022-2024). Comprobad si ese plan incluye objetivos
% específicos de gestión de personas / plantilla. La sección de
% Sostenibilidad de la Memoria Anual también incluye datos de plantilla
% y objetivos de RRHH.

[Descripción a completar con la información obtenida en la entrevista
y en los documentos públicos de CaixaBank (Memoria Anual, Plan Estratégico)]

\subsection{Carácter estratégico de la planificación de RRHH}

% Para que la planificación sea estratégica debe:
%   a) Estar integrada con el plan estratégico de negocio.
%   b) Tener un horizonte temporal de al menos 3 años.
%   c) Considerar escenarios de futuro (no solo extrapolar el pasado).
%   d) Involucrar a la alta dirección, no solo al departamento de RRHH.
%   e) Contar con indicadores de seguimiento (KPIs de plantilla).

\begin{table}[H]
  \centering
  \small
  \caption{Características de la planificación estratégica de RRHH en CaixaBank}
  \begin{tabularx}{\textwidth}{p{0.45\textwidth} p{0.12\textwidth} X}
    \toprule
    \textbf{Característica}                          & \textbf{¿Se cumple?} & \textbf{Observaciones} \\
    \midrule
    Integración con plan estratégico de negocio      & [Sí/No/Parcial] & \\
    Horizonte temporal $\geq$ 3 años                 & [Sí/No/Parcial] & \\
    Consideración de escenarios de futuro            & [Sí/No/Parcial] & \\
    Involucración de la alta dirección               & [Sí/No/Parcial] & \\
    KPIs de seguimiento definidos                    & [Sí/No/Parcial] & \\
    Planificación de la sucesión                     & [Sí/No/Parcial] & \\
    \bottomrule
  \end{tabularx}
\end{table}

\subsection{Impacto de la fusión con Bankia en la planificación de plantilla}

% La fusión CaixaBank-Bankia (2021) fue la mayor fusión bancaria de la
% historia de España. Implicó un ajuste de plantilla significativo (ERE
% 2021). Analizad cómo afectó a la planificación de RRHH: ¿fue reactiva
% o proactiva? ¿Cómo se gestionó la integración de plantillas?

[Descripción a completar]

\begin{notaequipo}
  Preguntad en la entrevista:
  \begin{enumerate}
    \item ¿Existe un plan de plantilla formal? ¿Quién lo elabora y con
          qué periodicidad?
    \item ¿Está vinculado al plan estratégico del grupo?
    \item ¿Cómo se gestionó la planificación de plantilla tras la fusión
          con Bankia?
    \item ¿Tienen planificación de sucesión para posiciones clave?
    \item ¿Qué KPIs de plantilla monitorizan (ratio empleado/oficina,
          rotación, cobertura de vacantes, etc.)?
  \end{enumerate}
\end{notaequipo}


% ================================================================
\section{Análisis de oferta y demanda de recursos humanos}
% ================================================================

% RECORDATORIO AL EQUIPO:
%
% ANÁLISIS DE DEMANDA (¿cuántos empleados necesitamos?):
%  Métodos cuantitativos:
%    - Análisis de series temporales (extrapolación de tendencias).
%    - Ratios de productividad (unidades producidas / empleado).
%    - Análisis de regresión.
%    - Modelos econométricos.
%  Métodos cualitativos:
%    - Técnica Delphi (consenso de expertos internos).
%    - Juicio directivo (managers estiman sus necesidades).
%    - Escenarios (planificación por escenarios).
%
% ANÁLISIS DE OFERTA (¿cuántos empleados tenemos / tendremos?):
%  Oferta interna (inventario de la plantilla actual):
%    - Inventario de habilidades / perfiles.
%    - Matrices de Markov / tablas de transición.
%    - Análisis de la rotación y del absentismo.
%    - Mapas de reemplazamiento / planes de sucesión.
%  Oferta externa (mercado laboral):
%    - Análisis del mercado de trabajo local / sectorial.
%    - Datos del INE, SEPE, convenios colectivos.
%    - Benchmarking con otras entidades financieras.

\subsection{Análisis de la demanda de RRHH}

La previsión de la demanda de personal permite a CaixaBank anticipar
cuántos y qué tipo de empleados necesitará en el futuro para alcanzar
sus objetivos de negocio.

\subsubsection{Métodos utilizados para la previsión de la demanda}

\begin{table}[H]
  \centering
  \small
  \caption{Métodos de previsión de la demanda de RRHH en CaixaBank}
  \begin{tabularx}{\textwidth}{p{0.30\textwidth} p{0.12\textwidth} X}
    \toprule
    \textbf{Método}                    & \textbf{¿Se usa?} & \textbf{Descripción / observaciones} \\
    \midrule
    Series temporales / tendencias     & [Sí/No/Parcial] & \\
    Ratios de productividad            & [Sí/No/Parcial] & \\
    Regresión / modelos cuantitativos  & [Sí/No/Parcial] & \\
    Técnica Delphi                     & [Sí/No/Parcial] & \\
    Juicio directivo                   & [Sí/No/Parcial] & \\
    Planificación por escenarios       & [Sí/No/Parcial] & \\
    \bottomrule
  \end{tabularx}
\end{table}

[Descripción detallada de los métodos más relevantes que usa CaixaBank]

\textbf{Métodos que serían convenientes aplicar:}

% TODO: Basándoos en el tamaño y complejidad de CaixaBank, ¿qué métodos
% más avanzados podrían añadir valor? Ejemplos:
%   - Modelos de regresión múltiple incorporando variables macroeconómicas
%     (volumen de crédito, tipos de interés, PIB).
%   - People Analytics con machine learning para predecir rotación.
%   - Planificación por escenarios ante disrupciones tecnológicas (IA).

\begin{itemize}
  \item [Método 1 — justificación]
  \item [Método 2 — justificación]
\end{itemize}

\subsection{Análisis de la oferta de RRHH}

\subsubsection{Oferta interna: inventario de la plantilla}

% TODO: ¿Tiene CaixaBank un sistema de inventario de habilidades?
% ¿Usan matrices de Markov o tablas de transición para modelar flujos
% internos? ¿Tienen mapas de reemplazamiento / planes de sucesión?

[Descripción a completar]

\subsubsection{Oferta externa: análisis del mercado de trabajo}

% TODO: ¿Cómo monitoriza CaixaBank el mercado laboral financiero?
% ¿Hacen benchmarking de salarios con competidores (Santander, BBVA)?
% ¿Colaboran con universidades para captar talento (prácticas, campus
% recruitment)?

[Descripción a completar]

\begin{notaequipo}
  Preguntad en la entrevista:
  \begin{enumerate}
    \item ¿Con qué métodos cuantifican sus necesidades futuras de personal?
    \item ¿Tienen un inventario actualizado de competencias de la plantilla?
    \item ¿Usan modelos matemáticos o estadísticos para prever la oferta
          interna (p.ej. matrices de transición / Markov)?
    \item ¿Cómo monitorizan el mercado laboral externo?
    \item ¿Cuál es su tasa de rotación actual? ¿La consideran alta o baja?
    \item ¿Han tenido dificultades para cubrir ciertos perfiles en los últimos
          años? ¿Cuáles? (Probablemente perfiles tecnológicos: Data Scientists,
          especialistas en ciberseguridad)
  \end{enumerate}
\end{notaequipo}


% ================================================================
\section{Equilibrio entre oferta y demanda de recursos humanos}
% ================================================================

% RECORDATORIO AL EQUIPO:
% Tras comparar oferta y demanda se pueden dar tres situaciones:
%
%  1. OFERTA = DEMANDA → Equilibrio. Ideal pero poco frecuente.
%     Acciones: mantenimiento, formación para cambios menores.
%
%  2. DEMANDA > OFERTA → Déficit de personal. Necesidad de incorporar.
%     Acciones: reclutamiento externo, promoción interna, horas extra,
%     trabajo temporal, outsourcing, boomerang employees, automatización.
%
%  3. OFERTA > DEMANDA → Superávit de personal.
%     Acciones: ERE, ERTE, jubilaciones anticipadas, congelación de
%     contrataciones, reducción de jornada, outplacement.
%     (Se desarrolla en la sección de procesos sustractivos)
%
% PARA ESTA SECCIÓN INDICAD:
%   - ¿Cuál es la situación de CaixaBank actualmente?
%   - ¿Ha variado en los últimos años (especialmente tras la fusión)?
%   - ¿Qué acciones concretas ha tomado para lograr el equilibrio?

\subsection{Situación actual de CaixaBank}

% TODO: Tras el ERE de 2021, la plantilla se ha ido ajustando. Pero la
% digitalización crea demanda de nuevos perfiles tecnológicos. Puede
% haber simultáneamente superávit en ciertos perfiles tradicionales
% (cajeros) y déficit en perfiles digitales.

[Descripción a completar]

\begin{table}[H]
  \centering
  \small
  \caption{Situación de equilibrio oferta-demanda en CaixaBank por tipo de perfil}
  \begin{tabularx}{\textwidth}{p{0.33\textwidth} p{0.22\textwidth} X}
    \toprule
    \textbf{Familia de puestos}    & \textbf{Situación} & \textbf{Acciones adoptadas} \\
    \midrule
    Perfiles bancarios tradicionales (oficina) & [Superávit/Déficit/Equilibrio] & \\
    Perfiles tecnológicos (datos, ciberseg., digital) & [Superávit/Déficit/Equilibrio] & \\
    Perfiles de gestión y dirección  & [Superávit/Déficit/Equilibrio] & \\
    Perfiles de cumplimiento y riesgo & [Superávit/Déficit/Equilibrio] & \\
    \bottomrule
  \end{tabularx}
\end{table}

\subsection{Acciones de ajuste ante déficit de personal}

\begin{itemize}
  \item \textbf{Reclutamiento externo:} [descripción]
  \item \textbf{Promoción interna:} [descripción]
  \item \textbf{Reskilling / upskilling:} [¿tienen programas de recualificación?]
  \item \textbf{Trabajo temporal o externalización:} [descripción]
  \item \textbf{Acuerdos con universidades / campus recruitment:} [descripción]
\end{itemize}

\subsection{Acciones de ajuste ante superávit de personal}

% Se desarrollará con mayor detalle en la sección de Procesos Sustractivos.
% Aquí haced un resumen de los mecanismos de reducción utilizados.

\begin{itemize}
  \item \textbf{ERE 2021 (fusión Bankia):} [número de afectados, condiciones]
  \item \textbf{Jubilaciones anticipadas:} [descripción]
  \item \textbf{Recolocación interna (movilidad geográfica / funcional):} [descripción]
  \item \textbf{Planes de formación para reconversión profesional:} [descripción]
\end{itemize}

\begin{notaequipo}
  Preguntad en la entrevista:
  \begin{enumerate}
    \item ¿En qué perfiles sienten más escasez de candidatos en el mercado?
    \item ¿Cómo afrontaron el exceso de plantilla tras la fusión con Bankia?
    \item ¿Han implementado programas de reskilling o upskilling para
          reconvertir perfiles tradicionales en perfiles digitales?
    \item ¿Cómo gestionan las situaciones de superávit temporal sin
          recurrir a despidos (ERTE, reducción de jornada...)?
    \item ¿Está la plantilla actual equilibrada o hay áreas con tensión?
  \end{enumerate}
\end{notaequipo}


% ================================================================
\section{Criterios de evaluación de la planificación y sistemas de información de RRHH}
% ================================================================

% RECORDATORIO AL EQUIPO:
%
% INDICADORES CLAVE DE RRHH (KPIs):
%   Indicadores de flujo de plantilla:
%     - Tasa de rotación (turnover rate): % de bajas sobre plantilla media.
%     - Tasa de absentismo: horas no trabajadas / horas contratadas.
%     - Tiempo medio de cobertura de vacantes (time-to-fill).
%     - Tasa de promoción interna.
%   Indicadores de calidad del ajuste:
%     - % de posiciones clave cubiertas por planes de sucesión.
%     - Tasa de retención en el primer año.
%   Indicadores de coste:
%     - Coste medio de contratación (cost-per-hire).
%     - ROI de formación.
%   Indicadores de previsión:
%     - Desviación entre plantilla planificada y real.
%
% SISTEMAS DE INFORMACIÓN DE RRHH (SIRH / HRIS):
%   Módulos típicos: administración de personal, gestión del tiempo,
%   selección (ATS), formación (LMS), evaluación del desempeño,
%   planificación de plantilla, People Analytics.
%   Sistemas del mercado: SAP SuccessFactors, Workday, Oracle HCM,
%   Cornerstone, Meta4 / NGA HR, A3 RRHH.

\subsection{Criterios e indicadores para evaluar la planificación de RRHH}

La eficacia de la planificación de RRHH se mide a través de indicadores
clave (KPIs) que permiten detectar desviaciones y tomar decisiones
correctoras.

\subsubsection{KPIs utilizados por CaixaBank}

\begin{table}[H]
  \centering
  \small
  \caption{Indicadores de planificación de RRHH en CaixaBank}
  \begin{tabularx}{\textwidth}{p{0.35\textwidth} p{0.15\textwidth} X}
    \toprule
    \textbf{Indicador}                    & \textbf{¿Se mide?} & \textbf{Valor / fuente} \\
    \midrule
    Tasa de rotación (turnover)           & [Sí/No] & [\% — Memoria Anual / entrevista] \\
    Tasa de absentismo                    & [Sí/No] & [\%] \\
    Tiempo medio de cobertura (TTF)       & [Sí/No] & [días] \\
    Tasa de promoción interna             & [Sí/No] & [\%] \\
    Coste medio de contratación           & [Sí/No] & [\euro{}] \\
    Desviación plantilla real vs. plan    & [Sí/No] & [\%] \\
    Cobertura de sucesión (puestos clave) & [Sí/No] & [\%] \\
    ROI de formación                      & [Sí/No] & [\%] \\
    \bottomrule
  \end{tabularx}
\end{table}

[Análisis de los indicadores más relevantes obtenidos en la entrevista]

\subsubsection{Propuestas de mejora en los criterios de evaluación}

% TODO: ¿Qué KPIs deberían medir y no miden? Por ejemplo:
%   - NPS de empleados (Employee Net Promoter Score).
%   - Índice de diversidad e inclusión.
%   - Indicadores de bienestar (wellbeing index).
%   - Métricas de People Analytics predictivo.

\begin{itemize}
  \item [Propuesta 1 — justificación]
  \item [Propuesta 2 — justificación]
\end{itemize}

\subsection{Sistemas de Información de Recursos Humanos (SIRH)}

\subsubsection{Sistema SIRH utilizado por CaixaBank}

% TODO: CaixaBank, dada su escala, es probable que use uno de los grandes
% sistemas enterprise (SAP SuccessFactors, Workday, Oracle HCM). En 2021,
% tras la fusión, habrán tenido que integrar los sistemas de ambas
% entidades. ¿Cuál es el sistema unificado resultante?

[Descripción a completar con la información de la entrevista]

\begin{table}[H]
  \centering
  \small
  \caption{Módulos del SIRH de CaixaBank}
  \begin{tabularx}{\textwidth}{p{0.42\textwidth} p{0.12\textwidth} X}
    \toprule
    \textbf{Módulo / Funcionalidad}          & \textbf{¿Activo?} & \textbf{Observaciones} \\
    \midrule
    Administración de personal / nómina       & [Sí/No] & \\
    Gestión del tiempo y asistencia           & [Sí/No] & \\
    Reclutamiento y selección (ATS)           & [Sí/No] & \\
    Gestión del desempeño                     & [Sí/No] & \\
    Formación y desarrollo (LMS)              & [Sí/No] & \\
    Planificación de plantilla                & [Sí/No] & \\
    People Analytics e informes               & [Sí/No] & \\
    Gestión de la compensación                & [Sí/No] & \\
    Portal del empleado (autoservicio)        & [Sí/No] & \\
    \bottomrule
  \end{tabularx}
\end{table}

\subsubsection{Propuestas de mejora del SIRH}

% TODO: Reflexionad sobre qué capacidades adicionales aportarían valor:
%   - Integrar People Analytics predictivo (predecir quién va a marcharse).
%   - Mejorar la experiencia del empleado en el portal de autoservicio.
%   - Integrar IA para el cribado de CV (con control del sesgo).
%   - Mejorar la interoperabilidad entre el SIRH y otros sistemas del banco.

[Descripción a completar]

\begin{notaequipo}
  Preguntad en la entrevista:
  \begin{enumerate}
    \item ¿Qué sistema SIRH utilizan? ¿Ha cambiado tras la fusión con Bankia?
    \item ¿Qué módulos tienen activos? ¿Qué módulos les gustaría activar?
    \item ¿Qué KPIs de plantilla monitorizan de forma habitual?
    \item ¿Tienen un cuadro de mando (dashboard) de RRHH? ¿Quién lo consulta?
    \item ¿Están satisfechos con la información que les proporciona el SIRH?
          ¿Qué mejorarían?
  \end{enumerate}
\end{notaequipo}


% ================================================================
\section{Procesos sustractivos: gestión del exceso de plantilla}
% ================================================================

% RECORDATORIO AL EQUIPO:
% Los procesos sustractivos son aquellos mediante los cuales la
% organización reduce su plantilla cuando la oferta de RRHH supera
% la demanda. Se clasifican en:
%
%  A) SUSPENSIÓN LABORAL TEMPORAL → ERTES
%     - El contrato se SUSPENDE temporalmente (no se extingue).
%     - El trabajador cobra prestación del SEPE + complemento empresa.
%     - Muy usado durante la pandemia COVID-19.
%     - Ventaja: mantiene la plantilla formada.
%
%  B) RUPTURA LABORAL VOLUNTARIA
%     1. Dimisión: el empleado abandona sin coste para la empresa.
%     2. Jubilación (ordinaria, anticipada o parcial).
%
%  C) RUPTURA LABORAL INVOLUNTARIA
%     1. Despido individual.
%     2. ERE — Expediente de Regulación de Empleo (despidos colectivos).
%     3. Downsizing: reducción sistemática del tamaño organizativo.
%     4. No renovación de contratos temporales.
%
%  D) RECOLOCACIÓN — OUTPLACEMENT
%     - Apoyo a los empleados que salen para que encuentren nuevo empleo.
%     - Mejora la imagen y reduce el trauma de la salida.

Cuando la oferta de RRHH supera la demanda, CaixaBank debe recurrir a
mecanismos de reducción o reajuste de plantilla. A continuación se
analizan los principales procesos sustractivos aplicados o aplicables.

\subsection{Suspensión laboral temporal: los ERTE}

% TODO: ¿Ha recurrido CaixaBank a ERTE? Principalmente durante la
% pandemia COVID-19. Si es así, ¿a cuántos empleados afectó? ¿Qué
% duración tuvo? ¿Qué complementos ofreció la empresa sobre el SEPE?

Un ERTE (Expediente de Regulación Temporal de Empleo) implica la
suspensión temporal del contrato de trabajo sin que se produzca su
extinción. El trabajador cobra una prestación por desempleo del SEPE,
que la empresa puede complementar.

\subsubsection{Aplicación en CaixaBank}

[Descripción a completar]

\textbf{Ventajas del ERTE para CaixaBank:}
\begin{itemize}
  \item [...]
\end{itemize}

\textbf{Inconvenientes / limitaciones:}
\begin{itemize}
  \item [...]
\end{itemize}

\subsection{Ruptura laboral voluntaria}

\subsubsection{Dimisión y rotación voluntaria}

% TODO: ¿Cuál es la tasa de dimisiones voluntarias en CaixaBank?
% ¿Realizan entrevistas de salida? ¿Qué factores llevan a los empleados
% a marcharse voluntariamente (mejores salarios en competencia, agotamiento...)?

[Descripción a completar]

\subsubsection{Jubilación (ordinaria y anticipada)}

% TODO: CaixaBank tiene una plantilla relativamente envejecida en
% ciertos segmentos. ¿Han implementado planes de jubilación anticipada?
% ¿En el ERE de 2021 se incluyeron medidas de jubilación anticipada?
% ¿Con qué condiciones (edad mínima, complemento sobre la pensión)?

[Descripción a completar]

\textbf{Ventajas de los planes de jubilación anticipada:}
\begin{itemize}
  \item Reducción no traumática de la plantilla sin despidos.
  \item Renovación generacional y entrada de perfiles más digitales.
  \item Menor conflicto sindical que los despidos colectivos.
\end{itemize}

\textbf{Inconvenientes:}
\begin{itemize}
  \item Pérdida de conocimiento experto y memoria institucional.
  \item Coste económico elevado para la empresa (complementos).
  \item Riesgo de que se jubilen empleados clave que no deberían marcharse.
\end{itemize}

\subsection{Ruptura laboral involuntaria: ERE y despidos}

\subsubsection{ERE 2021 — Fusión CaixaBank-Bankia}

% TODO: El ERE de 2021 fue el más grande de la historia bancaria española.
% Detalles clave a incluir:
%   - Número de afectados (aprox. 6.452 empleados según acuerdo sindical).
%   - Condiciones: indemnizaciones, bajas incentivadas, jubilaciones anticipadas.
%   - Sindicatos involucrados: UGT, CCOO, FINE.
%   - ¿Fue el ajuste suficiente o hay nuevos ajustes previstos?
%   - ¿Cómo afectó a la moral de la plantilla que se quedó?

En 2021, tras la fusión con Bankia, CaixaBank llevó a cabo el mayor
Expediente de Regulación de Empleo (ERE) de la banca española:

\begin{table}[H]
  \centering
  \small
  \caption{Datos del ERE CaixaBank 2021}
  \begin{tabularx}{\textwidth}{lX}
    \toprule
    \textbf{Parámetro}           & \textbf{Dato} \\
    \midrule
    Empleados afectados           & $\sim$6.452 (según acuerdo negociado) \\
    Cierre de oficinas            & $\sim$1.500 oficinas \\
    Tipo de salidas               & Bajas incentivadas, jubilaciones anticipadas, suspensiones \\
    Periodo de consultas          & [Fecha inicio -- Fecha acuerdo] \\
    Sindicatos firmantes          & [CCOO, UGT, FINE --- verificar] \\
    Medidas sociales de acompañamiento & [Outplacement, complementos, formación] \\
    \bottomrule
  \end{tabularx}
\end{table}

\textbf{Valoración del proceso:}

[Descripción a completar --- ¿fue percibido como justo? ¿Qué impacto
tuvo en el clima laboral? ¿Síndrome del superviviente?]

\subsubsection{Downsizing y reestructuración continua}

% TODO: Más allá del ERE de 2021, ¿hay un proceso de reducción continua
% (no renovación de contratos, digitalización que sustituye puestos)?
% La automatización de procesos bancarios reduce la necesidad de ciertos
% perfiles de back-office.

[Descripción a completar]

\subsection{Recolocación: el \textit{outplacement}}

\subsubsection{Servicios de \textit{outplacement} en el ERE de CaixaBank}

% TODO: ¿Ofreció CaixaBank servicios de outplacement a los empleados
% que salieron? ¿Lo gestionaron internamente o con consultora externa
% (Lee Hecht Harrison, Randstad RiseSmart, Adecco)?
% ¿Cuántos empleados se beneficiaron? ¿Con qué tasa de recolocación?

El outplacement es un conjunto de servicios de apoyo a los empleados
que abandonan la empresa, orientados a facilitar su reincorporación
al mercado laboral:

\begin{itemize}
  \item Orientación y coaching profesional.
  \item Actualización y elaboración del currículum.
  \item Preparación para entrevistas de trabajo.
  \item Identificación de oportunidades en el mercado.
  \item Apoyo psicológico durante la transición.
\end{itemize}

[Descripción de cómo se aplicó en CaixaBank]

\textbf{Ventajas del outplacement para CaixaBank:}
\begin{itemize}
  \item Mejora la imagen como empleador responsable (\textit{employer branding}).
  \item Reduce el conflicto y los litigios laborales.
  \item Mitiga el impacto emocional en los empleados que salen, lo que
        también reduce el síndrome del superviviente en los que se quedan.
\end{itemize}

\begin{notaequipo}
  Preguntad en la entrevista:
  \begin{enumerate}
    \item ¿Han recurrido a ERTE en algún momento (COVID, fusión)?
    \item ¿Qué condiciones se ofrecieron en el ERE de 2021?
          ¿Hubo jubilaciones anticipadas? ¿A partir de qué edad?
    \item ¿Se ofreció outplacement a los empleados afectados por el ERE?
          ¿Quién lo gestionó? ¿Cuáles fueron los resultados?
    \item ¿Cómo afectó el ERE al clima laboral de los que se quedaron?
          ¿Tomaron medidas para gestionar ese impacto?
    \item ¿Realizan entrevistas de salida cuando un empleado dimite?
          ¿Qué hacen con esa información?
  \end{enumerate}
\end{notaequipo}


% ================================================================
\section{Alternativas a la reducción de plantilla e incorporación a la empresa}
% ================================================================

% RECORDATORIO AL EQUIPO:
%
% ALTERNATIVAS A LA REDUCCIÓN (cuando hay superávit):
%   1. Reducción de la jornada laboral (trabajo a tiempo parcial).
%   2. Reparto del trabajo (job sharing / worksharing).
%   3. Reducción salarial temporal (sacrificio compartido).
%   4. Congelación de contrataciones y de salarios.
%   5. Formación y recualificación (reskilling).
%   6. Movilidad funcional: asignar empleados a otras áreas con demanda.
%   7. Movilidad geográfica: trasladar empleados a zonas con necesidades.
%   8. Excedencias voluntarias y licencias sin sueldo.
%   9. Reducción de horas extra.
%  10. Externalización de funciones no estratégicas (outsourcing).
%
% ALTERNATIVAS A LA INCORPORACIÓN (cuando hay déficit):
%   1. Promoción interna (movilidad vertical).
%   2. Formación acelerada para cubrir el gap de habilidades.
%   3. Horas extra (si el déficit es temporal y limitado).
%   4. Trabajo temporal o contratos de proyecto.
%   5. Externalización (BPO, outsourcing, contratas).
%   6. Automatización (RPA, IA) para reducir la necesidad de personal.
%   7. Programas de boomerang employees (recontratación de exempleados).

\subsection{Alternativas a la reducción de plantilla}

Antes de recurrir a medidas de ruptura laboral, CaixaBank puede ---y
debe--- explorar alternativas menos traumáticas para ajustar la
plantilla preservando el empleo y el clima organizacional.

\begin{table}[H]
  \centering
  \small
  \caption{Alternativas a la reducción de plantilla en CaixaBank}
  \begin{tabularx}{\textwidth}{p{0.37\textwidth} p{0.12\textwidth} X}
    \toprule
    \textbf{Alternativa}                        & \textbf{¿Se aplica?} & \textbf{Observaciones} \\
    \midrule
    Reducción de jornada / trabajo parcial       & [Sí/No] & \\
    Congelación de contrataciones                & [Sí/No] & \\
    Reskilling / reconversión de perfiles        & [Sí/No] & \\
    Movilidad funcional interna                  & [Sí/No] & \\
    Movilidad geográfica                         & [Sí/No] & \\
    Excedencias voluntarias incentivadas         & [Sí/No] & \\
    Outplacement previo al despido               & [Sí/No] & \\
    Reducción de horas extra / subcontratación   & [Sí/No] & \\
    \bottomrule
  \end{tabularx}
\end{table}

\subsubsection{Reskilling y recualificación profesional}

% TODO: Este es probablemente el aspecto más relevante para CaixaBank.
% La digitalización bancaria hace que muchos perfiles de red de oficinas
% queden obsoletos. ¿Tiene CaixaBank un programa de reskilling para
% reconvertir a estos empleados en perfiles digitales?
% Mencionad si existe un programa concreto (p.ej. ``Future Skills'' o similar).

[Descripción a completar]

\textbf{Ventajas:}
\begin{itemize}
  \item Preserva el empleo y el conocimiento del negocio.
  \item Mejor imagen como empresa responsable.
  \item Menor coste que despedir y contratar externamente.
\end{itemize}

\textbf{Inconvenientes:}
\begin{itemize}
  \item No todos los empleados tienen el potencial o la motivación para
        reconvertirse en perfiles muy técnicos.
  \item Requiere inversión elevada en formación.
  \item El proceso es lento; no resuelve necesidades inmediatas de
        perfiles digitales especializados.
\end{itemize}

\subsubsection{Movilidad interna (funcional y geográfica)}

% TODO: ¿Tiene CaixaBank un sistema formal de movilidad interna?
% ¿Los empleados pueden postularse a vacantes internas? ¿Hay un portal
% de oportunidades internas en el SIRH? ¿Con qué frecuencia se producen
% traslados voluntarios o forzosos?

[Descripción a completar]

\subsection{Alternativas a la incorporación de nueva plantilla}

Cuando existe déficit de personal, CaixaBank no siempre recurre de
inmediato a la contratación externa. Las alternativas disponibles
incluyen:

\subsubsection{Promoción interna}

% TODO: ¿Cuál es la política de CaixaBank respecto a la promoción
% interna? ¿Se publican primero las vacantes internamente antes de
% salir al mercado? ¿Tienen planes de carrera definidos?

[Descripción a completar]

\subsubsection{Automatización y digitalización}

% TODO: La RPA y la IA están transformando el sector bancario. CaixaBank
% ha invertido fuertemente en tecnología. ¿Hay procesos que antes
% requerían empleados y ahora están automatizados? ¿Esto ha reducido
% la necesidad de incorporaciones en ciertas áreas?

[Descripción a completar]

\subsubsection{Externalización (\textit{outsourcing})}

% TODO: ¿Qué funciones tiene CaixaBank externalizadas?
% ¿IT, call centers, servicios de back-office, seguridad, limpieza?
% ¿Han seguido la tendencia de insourcing en algún área?

[Descripción a completar]

\subsubsection{Trabajo temporal y contratos de proyecto}

[Descripción a completar]

\begin{notaequipo}
  Preguntad en la entrevista:
  \begin{enumerate}
    \item ¿Tienen algún programa de reskilling o upskilling para reconvertir
          perfiles tradicionales? ¿En qué consiste? ¿Cuántos empleados se
          han beneficiado?
    \item ¿Existe un portal de vacantes internas para la movilidad interna?
    \item ¿Cómo deciden entre cubrir una vacante internamente o acudir
          al mercado externo?
    \item ¿Qué procesos han automatizado recientemente con RPA o IA?
          ¿Ha supuesto una reducción de la plantilla en esas áreas?
    \item ¿Qué funciones tienen externalizadas? ¿Han internalizado alguna
          función en los últimos años?
  \end{enumerate}
\end{notaequipo}
